\documentclass[11pt,a4paper]{article}
\usepackage[utf8]{inputenc}
\usepackage[T1]{fontenc}
\usepackage[brazil]{babel}
\usepackage{lmodern}
\usepackage[margin=2cm]{geometry}
\usepackage{amsmath,amssymb}
\usepackage{booktabs}
\usepackage{graphicx}
\usepackage{hyperref}
\usepackage{xcolor}
\usepackage{titlesec}
\usepackage{tabularx}
\usepackage{enumitem}
\usepackage{float}
\usepackage{array}
\renewcommand{\tabularxcolumn}[1]{>{\raggedright\arraybackslash}p{#1}}
\usepackage{newunicodechar}
\usepackage[most]{tcolorbox}
\usepackage{etoolbox}

\hypersetup{colorlinks=true, linkcolor=blue!70!black, citecolor=blue!70!black, urlcolor=blue!70!black}

\titleformat{\section}{\large\bfseries}{\thesection.}{0.5em}{}
\titleformat{\subsection}{\normalsize\bfseries}{\thesubsection.}{0.5em}{}
\titlespacing*{\section}{0pt}{1.2ex plus 0.5ex minus .2ex}{0.6ex}
\titlespacing*{\subsection}{0pt}{1.0ex plus 0.4ex minus .2ex}{0.4ex}
\setlength{\parskip}{0.45em}
\setlength{\parindent}{0pt}
\setlist{itemsep=0.15em,topsep=0.3em}

% --- símbolos ---
\newunicodechar{π}{\ensuremath{\pi}}
\newunicodechar{τ}{\ensuremath{\tau}}
\newunicodechar{φ}{\ensuremath{\varphi}}
\newunicodechar{ψ}{\ensuremath{\psi}}
\newunicodechar{Σ}{\ensuremath{\Sigma}}
\newunicodechar{→}{\ensuremath{\rightarrow}}
\newunicodechar{⇒}{\ensuremath{\Rightarrow}}
\newunicodechar{∈}{\ensuremath{\in}}
\newunicodechar{−}{\ensuremath{-}}
\newunicodechar{≈}{\ensuremath{\approx}}
\newunicodechar{≠}{\ensuremath{\neq}}
\newunicodechar{≡}{\ensuremath{\equiv}}
\newunicodechar{≤}{\ensuremath{\leq}}
\newunicodechar{≥}{\ensuremath{\geq}}
\newunicodechar{′}{\ensuremath{'}}
\newunicodechar{‑}{-}
\newunicodechar{✔}{\ensuremath{\checkmark}}
\newunicodechar{✘}{\ensuremath{\times}}

% --- cores e caixas ---
\definecolor{cC}{HTML}{7A4FA3}
\definecolor{cS}{HTML}{2F7D74}
\definecolor{cAcc}{HTML}{2F5D8A}
\newcommand{\metrictag}[2]{{\setlength{\fboxsep}{1.5pt}%
  \ifstrequal{#1}{c}{\colorbox{cC}{\textcolor{white}{\scriptsize\textbf{#2}}}}{%
  \ifstrequal{#1}{s}{\colorbox{cS}{\textcolor{white}{\scriptsize\textbf{#2}}}}{%
  \colorbox{black!60}{\textcolor{white}{\scriptsize\textbf{#2}}}}}}}
\newcommand{\pagechip}[1]{{\setlength{\fboxsep}{1.5pt}\colorbox{cAcc}{\textcolor{white}{\scriptsize\texttt{#1}}}}}
\newcommand{\readtag}[1]{\texttt{\small #1}}
\newcommand{\partbar}[1]{\par\vspace{1.2em}\noindent\colorbox{black!80}{\parbox{\dimexpr\linewidth-2\fboxsep}{\color{white}\bfseries\sffamily\MakeUppercase{#1}}}\par\vspace{0.4em}}
\newtcolorbox{citacard}[1]{breakable, enhanced, colback=white, colframe=black!25, boxrule=0.5pt, arc=2pt,
  left=6pt, right=6pt, top=4pt, bottom=4pt, title={\small #1}, coltitle=black, colbacktitle=black!6, fonttitle=\normalfont}
\newenvironment{origquote}{\par\itshape\small}{\par}
\newenvironment{tradquote}{\par\small\textsc{Tradução:}\ }{\par}
\newcommand{\explab}[1]{\par\smallskip{\small\color{cAcc}\textbf{\MakeUppercase{#1}}}\par}
\colorlet{notebg}{blue!4}\colorlet{notefg}{cAcc}
\colorlet{okbg}{green!5}\colorlet{okfg}{green!45!black}
\colorlet{warnbg}{orange!8}\colorlet{warnfg}{orange!80!black}
\newtcolorbox{infobox}[2][note]{breakable, enhanced, boxrule=0pt, arc=2pt, left=6pt, right=6pt, top=4pt, bottom=4pt,
  colback=#1bg, borderline west={3pt}{0pt}{#1fg}, frame hidden,
  before upper={\ifstrempty{#2}{}{\textbf{#2}\par}}}
\colorlet{pairc}{cC}\colorlet{pairs}{cS}
\newenvironment{paircol}[2]{\begin{tcolorbox}[width=0.485\linewidth, nobeforeafter, box align=top, boxrule=0.5pt,
  colframe=black!25, colback=white, borderline north={3pt}{0pt}{pair#1}, left=5pt, right=5pt,
  before upper={{\color{pair#1}\small\textbf{\MakeUppercase{#2}}}\par}]}{\end{tcolorbox}}

\newunicodechar{ℓ}{\ensuremath{\ell}}
\newunicodechar{⊙}{\ensuremath{\odot}}
\newunicodechar{α}{\ensuremath{\alpha}}
\newunicodechar{β}{\ensuremath{\beta}}

\begin{document}

\begin{center}
    {\LARGE \textbf{Fichamento: sensibilidade de explicações em séries temporais}} \\[0.35em]
    {\large FOULADGAR, N.; ALIREZAIE, M.; FRÄMLING, K. \textit{Metrics and Evaluations of Time Series Explanations: An Application in Affect Computing}. IEEE Access, v.~10, p.~23995--24009, 2022} \\[0.3em]
    \textbf{Tema:} métricas de sensibilidade (Max/Avg) adaptadas a séries temporais e o efeito da escala dos escores \\[0.2em]
    \textbf{Aluno:} Arthur Santos de Oliveira (RA: 156297) \\[0.2em]
    \textbf{Atualizado em:} 07/11/2026 \\[0.2em]
    \textbf{Link:} \url{https://doi.org/10.1109/ACCESS.2022.3155115} \\[0.2em]
    \textbf{Leitura:} íntegra (versão publicada, 15 páginas; PDF baixado pelo aluno) \\[0.2em]
    \textbf{Repositório:} \url{https://github.com/Arthur-so/shap-ecg-arrhythmia}
\end{center}
\vspace{-0.5em}
\hrule
\vspace{0.8em}

\textbf{Como ler.} \texttt{[p. N]} indica a página do PDF (página da revista $= 23994 + N$). As citações estão no original em inglês, cada uma seguida de uma tradução livre e de uma explicação.

\begin{infobox}[note]{Resumo}
\begin{itemize}
\item \textbf{Objetivo:} propor métricas de sensibilidade de explicações para séries temporais e avaliá-las em computação afetiva.
\item \textbf{Métricas:} sensibilidade máxima e média, de curto prazo (série perturbada por ruído) e de longo prazo (vizinhos reais do treino com a mesma classe): MSS, ASS, MLS e ALS.
\item \textbf{Dados e modelos:} sinais fisiológicos multivariados (WESAD e MAHNOB-HCI, janelas de 30 instantes) e uma CNN (CN-Waterfall); explicações por LIME, gradientes integrados (IG) e SmoothGrad (SG).
\item \textbf{Achado principal:} IG e LIME parecem menos sensíveis que o SG em todas as métricas, \textbf{possivelmente só porque os seus escores têm escala menor}; a sensibilidade também varia com o tamanho da janela, com o treino do modelo e com os hiperparâmetros dos métodos.
\end{itemize}
\end{infobox}

\section{O que o trabalho faz}
As métricas seguem a Max-Sensitivity e a Avg-Sensitivity: a maior e a média das distâncias entre a explicação da série e as de séries próximas com a mesma predição [p.~5]. O ruído de curto prazo é gaussiano, e os vizinhos de longo prazo são séries reais do treino [p.~4--5]. Os autores variam a janela e a sobreposição dos dados, o treino do modelo (épocas) e os hiperparâmetros de cada método (passos do IG, desvio do ruído do SG, vizinhança do LIME) [p.~7--13].

\section{Citações}

\begin{citacard}{\pagechip{p. 1} Resumo: o efeito da escala}
\begin{origquote}“The results reveal that (i) IG and LIME provide a lower sensitivity scale than SG in all the metrics and settings, potentially due to the lower scale of important scores generated by IG and LIME [...]”\end{origquote}
\begin{tradquote}“Os resultados mostram que (i) IG e LIME têm sensibilidade em escala menor que o SG em todas as métricas e configurações, possivelmente por causa da escala menor dos escores de importância gerados por IG e LIME [...]”\end{tradquote}
\explab{Explicação}
Se a distância entre mapas não é normalizada, um método que gera escores menores parece mais robusto sem ser. É o motivo de o TCC normalizar os vetores das 36 regiões para norma 1 antes de calcular qualquer distância (piso, sensibilidades, Random Logit, eMPRT).
\end{citacard}

\begin{citacard}{\pagechip{p. 7} A escala dos escores}
\begin{origquote}“We also infer that the scale of importance scores in IG and LIME are much lower than in SG (nearly close to 0). Such scaling could result in lower sensitivities in the former models than in the latter.”\end{origquote}
\begin{tradquote}“Inferimos também que a escala dos escores de importância de IG e LIME é muito menor que a do SG (quase zero). Essa escala pode resultar em sensibilidades menores nos primeiros do que no último.”\end{tradquote}
\explab{Explicação}
Os autores não testam a hipótese com distâncias normalizadas; a frase deve ser citada como indício (``possivelmente''), e não como demonstração.
\end{citacard}

\begin{citacard}{\pagechip{p. 14} Conclusão}
\begin{origquote}“In summary, we found that (i) IG and LIME provide lower scales of sensitivity than SG in all the metrics and settings. We referred the result to the lower scale of important scores generated by the former models; (ii) the window size of the series plays a role in the sensitivities’ variation in the XAI models.”\end{origquote}
\begin{tradquote}“Em resumo, verificamos que (i) IG e LIME têm sensibilidades em escala menor que o SG em todas as métricas e configurações, o que atribuímos à escala menor dos escores gerados por esses métodos; (ii) o tamanho da janela da série influencia a variação das sensibilidades.”\end{tradquote}
\explab{Explicação}
O segundo ponto também interessa ao TCC: a sensibilidade depende da unidade em que a série é avaliada, o que reforça calcular todas as métricas na mesma unidade (as 36 regiões).
\end{citacard}

\section{Relação com o TCC}
É um precedente das métricas de sensibilidade em séries temporais e o apoio para normalizar os mapas antes de compará-los. Diferenças: sinais de computação afetiva, janelas curtas (30 instantes), métodos sem GradientSHAP, nenhuma leitura em relação à variação própria do método (o piso do TCC) e nenhuma comparação por classe.

\section{Como citar no TCC}
\begin{itemize}
\item ``Em séries temporais, Fouladgar et al.\ (2022) observam que métodos cujos escores têm escala menor aparecem como menos sensíveis, possivelmente só por causa da escala; por isso, os mapas foram normalizados antes do cálculo das distâncias.''
\end{itemize}

\end{document}
